\section{算力资源：让计算算得快、存得下、传得动}
\label{sec:hw}

\begin{guide}
\begin{minipage}[c]{0.66\linewidth}
\noindent 上一节的估算给出了三项硬要求：高算力、大容量高带宽的存储，以及高速互联。硬件如何满足它们？计算芯片提供算力，HBM 等存储提供容量与带宽，NVLink、InfiniBand 等互联技术把多卡多机连成整体，驱动与计算软件栈再把这些能力开放给上层软件。不过要记住，硬件只决定性能的上界，能否接近上界，还要看软件能否用好它。
\end{minipage}\hfill
\begin{minipage}[c]{0.3\linewidth}\centering\locator{4}\end{minipage}
\end{guide}

\subsection{计算芯片：CPU、GPU 与 NPU/ASIC}
\label{subsec:chips}

先看大模型到底在算什么。无论是注意力中的投影、注意力分数与加权求和，还是前馈网络中的两次线性变换，Transformer 的绝大部分计算都可以归结为矩阵乘法。矩阵乘法由海量彼此独立的乘加运算组成，谁先算、谁后算都不影响结果。因此，决定速度的不是单个运算有多快，而是\emph{同一时刻能有多少运算并行进行}——大模型需要的是高并行吞吐，而不是单个任务的低延迟。

哪种芯片能提供这样的吞吐？不同芯片在“通用性”与“专用效率”之间做出了不同取舍（图\,\ref{fig:chips}）。
\begin{itemize}
  \item \term{CPU}（\en{Central Processing Unit}，中央处理器）核心数较少，但每个核心具有复杂的分支预测和多级缓存，擅长控制流密集、需要低延迟响应的通用任务。在大模型系统中，CPU 通常负责操作系统、数据预处理、请求编排、内存管理，以及向加速器下发计算任务。
  \item \term{GPU}（\en{Graphics Processing Unit}，图形处理器）集成了成千上万个并行执行单元，并配有专门的矩阵乘加单元（如 NVIDIA 的 Tensor Core）和高带宽显存，适合训练和推理中的大规模线性代数运算\cite{hwu2026pmpp}。
  \item \term{NPU}（\en{Neural Processing Unit}，神经网络处理器）泛指为神经网络计算专门优化的处理器；\term{ASIC}（\en{Application-Specific Integrated Circuit}，专用集成电路）则泛指为特定用途定制设计的芯片。二者并非互斥的分类——前者按\emph{用途}命名，后者按\emph{实现方式}命名，一款 NPU 完全可以用 ASIC 方式实现。Google TPU、华为昇腾 NPU 和寒武纪 MLU 都属于这一类。
\end{itemize}

\begin{figure}[htbp]
\centering
\begin{tikzpicture}[fig, x=1mm, y=1mm]
  % CPU：少量复杂核心 + 大缓存
  \begin{scope}[shift={(-48,0)}]
    \draw[figline, fill=white] (-17,-14) rectangle (17,14);
    \foreach \x/\y in {-8/7,8/7,-8/-1,8/-1}
      \draw[fill=figfill] (\x-6,\y-3.5) rectangle (\x+6,\y+3.5);
    \foreach \x/\y in {-8/7,8/7,-8/-1,8/-1} \node[font=\figfont\scriptsize] at (\x,\y) {核心};
    \draw[fill=white, draw=rulegray] (-14,-12) rectangle (14,-6);
    \node[font=\figfont\scriptsize] at (0,-9) {大容量缓存};
    \node[anchor=north] at (0,-16) {CPU};
    \node[note, anchor=north] at (0,-21) {少量复杂核心};
  \end{scope}
  % GPU：大量简单核心
  \begin{scope}[shift={(0,0)}]
    \draw[figline, fill=white] (-17,-14) rectangle (17,14);
    \foreach \i in {0,...,7} \foreach \j in {0,...,6}
      \draw[fill=figfill, line width=0.3pt] (-15+\i*3.8,-12+\j*3.8) rectangle ++(3,3);
    \node[anchor=north] at (0,-16) {GPU};
    \node[note, anchor=north] at (0,-21) {大量简单并行单元};
  \end{scope}
  % NPU：矩阵运算阵列
  \begin{scope}[shift={(48,0)}]
    \draw[figline, fill=white] (-17,-14) rectangle (17,14);
    \foreach \i in {0,...,5} \foreach \j in {0,...,5}
      \draw[fill=accentpale, draw=accent, line width=0.3pt] (-12+\i*4.2,-12+\j*4.2) rectangle ++(3.4,3.4);
    \foreach \j in {0,...,5} \draw[->, accent, line width=0.4pt] (-16,-10.3+\j*4.2) -- (-12.4,-10.3+\j*4.2);
    \node[anchor=north] at (0,-16) {NPU / ASIC};
    \node[note, anchor=north] at (0,-21) {专用矩阵运算阵列};
  \end{scope}
  % 光谱箭头
  \draw[<->, line width=0.7pt] (-65,-30) -- (65,-30);
  \node[note, anchor=north west] at (-65,-31) {通用性、灵活性更强};
  \node[note, anchor=north east] at (65,-31) {专用计算效率更高};
\end{tikzpicture}
\caption{三类计算芯片的结构取舍（示意）。从左到右，芯片的通用性逐渐降低，矩阵运算的吞吐与能效逐渐提高}
\label{fig:chips}
\end{figure}


实际系统中，CPU 与加速器是\emph{协同}关系而非替代关系：CPU 负责发起和调度，GPU 或 NPU 负责执行吞吐密集的张量运算。表\,\ref{tab:chips} 对几类芯片做了概括。

\begin{table}[htbp]
\centering
\caption{计算芯片的适用性比较}
\label{tab:chips}
\small
\begin{tabularx}{\linewidth}{@{}l Y Y Y@{}}
\toprule
\heiti 类型 & \heiti 典型优点 & \heiti 典型限制 & \heiti 在大模型系统中的职责 \\
\midrule
CPU & 控制灵活，软件生态成熟 & 张量吞吐远低于专用加速器 & 预处理、调度、系统服务、小模型推理 \\
GPU & 并行吞吐高，深度学习生态成熟 & 显存容量有限，多卡互联成本高 & 训练、微调、高吞吐推理 \\
NPU / ASIC & 针对矩阵和低精度运算优化，能效高 & 框架、算子与驱动的兼容性因平台而异 & 云端或边缘侧的专用 AI 计算 \\
FPGA & 可重构，可定制数据通路 & 开发难度大，生态不如 GPU & 特定低延迟或定制化场景 \\
\bottomrule
\end{tabularx}
\end{table}

专用化换来了更高的吞吐与能效，代价又是什么？一是通用性下降，二是软件生态的绑定加深。这也提醒我们，选择芯片时不能只比较“TFLOPS”或“TOPS”一个数字：不同精度（FP16、BF16、FP8、INT8）、是否计入稀疏加速，测得的峰值差别可达数倍。更可靠的做法是同时比较目标模型的支持情况、低精度格式、板载存储容量、编译器与算子库的成熟度、框架适配、集群通信能力，以及采购和运维成本。

\subsection{存储层次：为什么 HBM 如此重要}
\label{subsec:hbm}

第\,\ref{subsec:three-quantities}\,小节留下了一个结论：解码时，每生成一个 token 都要把全部权重从显存读一遍，速度上限由显存带宽决定。训练时情况更紧张——权重、梯度、优化器状态和激活值都要驻留在设备上。于是存储系统要同时回答两个问题：\emph{装不装得下}，以及\emph{送不送得及}。

单一一种存储无法同时做到又大、又快、又便宜，计算机的办法是\emph{分层}（图\,\ref{fig:memory}）。越靠近计算单元，存储越快、越小、越贵：片上 SRAM（寄存器、共享内存、缓存）带宽最高但容量只有数十 MB；板载显存容量为数十 GB；主机内存（DRAM）可达 TB 级，但 GPU 访问它要经过 PCIe，带宽比显存低一个数量级以上；SSD 用于存放模型文件和冷数据。

\begin{figure}[htbp]
\centering
\begin{tikzpicture}[fig, x=1mm, y=1mm]
  % 金字塔：每层梯形，顶部最窄
  \def\h{9}
  \foreach \k/\name/\fillc/\drawc in {0/片上 SRAM/accentpale/accent, 1/HBM 显存/figfill/figline, 2/主机内存 DRAM/white/figline, 3/SSD 存储/white/rulegray} {
    \pgfmathsetmacro{\yt}{-\k*\h}
    \pgfmathsetmacro{\yb}{-(\k+1)*\h}
    \pgfmathsetmacro{\wt}{9+\k*7}
    \pgfmathsetmacro{\wb}{9+(\k+1)*7}
    \draw[draw=\drawc, fill=\fillc] (-\wt,\yt) -- (\wt,\yt) -- (\wb,\yb) -- (-\wb,\yb) -- cycle;
    \node[font=\figfont\small] at (0,{(\yt+\yb)/2}) {\name};
  }
  % 表头
  \node[note, anchor=west, font=\figfont\footnotesize\bfseries] at (42,4) {典型容量};
  \node[note, anchor=west, font=\figfont\footnotesize\bfseries] at (63,4) {对 GPU 的带宽};
  \foreach \k/\c/\b in {0/{数十 MB}/{数十 TB/s}, 1/{数十 GB}/{约 3 TB/s}, 2/{数百 GB～数 TB}/{经 PCIe 约 64 GB/s}, 3/{数 TB 以上}/{数 GB/s～十余 GB/s}} {
    \pgfmathsetmacro{\yc}{-(\k+0.5)*\h}
    \pgfmathsetmacro{\xs}{9+(\k+0.5)*7+2}
    \node[note, anchor=west] at (42,\yc) {\c};
    \node[note, anchor=west] at (63,\yc) {\b};
    \draw[rulegray, line width=0.3pt, dotted] (\xs,\yc) -- (41,\yc);
  }
  % 左侧趋势箭头
  \draw[->, line width=0.7pt] (-42,-36) -- (-42,0) node[pos=0.5, left, note, text width=8mm, align=right] {更快\\更小\\更贵};
\end{tikzpicture}
\caption{加速器视角下的存储层次（数量级示意，以 H100 级数据中心 GPU 为参照）。模型权重和 KV 缓存主要驻留在 HBM 中，计算时被反复读入片上 SRAM}
\label{fig:memory}
\end{figure}


在这个层次中，显存这一级对大模型最为关键，而\term{高带宽内存}（\en{High Bandwidth Memory}，HBM）是当前数据中心加速器的主流显存方案。它把多层 DRAM 芯片垂直堆叠，用\term{硅通孔}（\en{Through-Silicon Via}，TSV）实现层间互联，再通过很宽的接口（每个堆栈 1024 位）与计算芯片封装在一起。接口宽、距离短，使 HBM 能在有限的功耗和面积内提供远高于 GDDR 等传统显存的带宽。例如 H100 SXM 配备 \SI{80}{GB} HBM3，带宽约 \SI{3.35}{TB/s}\cite{nvidia2023h100}。

不过，HBM 解决的主要是“送不送得及”，而不是“装不装得下”。它价格昂贵，容量依然有限；运行大模型也并非必须使用 HBM——消费级显卡的 GDDR 显存同样可以运行规模适中或经过量化的模型。不妨把三个指标的分工记在一起：\emph{容量决定模型和状态能否装下，带宽决定数据能否及时送达，延迟影响细粒度访问的效率}。

\begin{pitfall}
“7B 模型 4 位量化后仅权重就约 \SI{3.5}{GB}”不等于“有 \SI{4}{GB} 显存就能运行”。实际运行还需要存放量化缩放系数等元数据、KV 缓存、运行时缓冲区和临时张量，并存在内存碎片。估算时应在权重之外预留足够余量。
\end{pitfall}

\subsection{互联技术：PCIe、NVLink 与 InfiniBand}
\label{subsec:interconnect}

一张卡总有装不下、算不完的时候。例\,\ref{ex:gpus} 表明，大模型训练需要成百上千张卡协同工作；即使只做推理，较大的模型也常常要切分到多张卡上。一旦切分，卡与卡之间就要不断交换梯度、激活值和中间结果。此时计算单元再快，如果数据迟迟不到，也只能空等。

讨论互联时，最有用的问题不是“哪种技术更快”，而是“\emph{数据要从哪里传到哪里}”。不同层级的链路各有分工（图\,\ref{fig:interconnect}）：
\begin{itemize}
  \item \term{PCIe}（\en{PCI Express}）是通用的设备互联总线，连接 CPU 与 GPU、网卡、NVMe 硬盘等外设。PCIe 5.0 x16 的单向带宽约 \SI{64}{GB/s}。
  \item \term{NVLink} 是 NVIDIA 的 GPU 高速互联技术，用于同一服务器内 GPU 与 GPU（或与 Grace CPU）之间的直接通信，带宽远高于 PCIe（H100 每卡 NVLink 总带宽约 \SI{900}{GB/s}\cite{nvidia2023h100}）。配合 NVSwitch 交换芯片，可以让一台服务器内的 8 张 GPU 两两全速互联。
  \item \term{InfiniBand}（IB）是面向高性能计算集群的\emph{跨节点}网络，单端口带宽可达 \SI{400}{Gb/s}（NDR）。它原生支持 \term{RDMA}（\en{Remote Direct Memory Access}，远程直接内存访问）：数据可以直接写入远端机器的内存，而不必经过双方 CPU 的多次拷贝和内核协议栈处理。以太网借助 \term{RoCE}（\en{RDMA over Converged Ethernet}）也能实现 RDMA，成本较低，但在拥塞控制和延迟稳定性上通常需要更细致的调优。
\end{itemize}

\begin{figure}[htbp]
\centering
\begin{tikzpicture}[fig, x=1mm, y=1mm]
  \foreach \s/\nm in {-1/A, 1/B} {
    \begin{scope}[shift={(\s*36,0)}]
      \node[grp, minimum width=62mm, minimum height=44mm] (box\nm) at (0,0) {};
      \node[anchor=north west, font=\figfont\footnotesize\bfseries] at (box\nm.north west) {服务器 \nm};
      % NVSwitch
      \node[blka, minimum width=44mm, minimum height=6mm, inner sep=1pt, font=\figfont\footnotesize] (sw\nm) at (0,9) {NVSwitch};
      % GPU ×4（示意，实际常为 8）
      \foreach \i in {0,...,3} {
        \node[blk, minimum width=9mm, minimum height=7mm, inner sep=1pt, font=\figfont\footnotesize] (g\nm\i) at (-16.5+\i*11,-1) {GPU};
        \draw[accent, line width=1.4pt] (g\nm\i.north) -- (g\nm\i.north |- sw\nm.south);
      }
      \node[font=\figfont\scriptsize, anchor=west] at (22,-1) {×8};
      % PCIe 总线
      \draw[figline, line width=0.6pt] (-24,-10) -- (24,-10);
      \foreach \i in {0,...,3} \draw[figline] (g\nm\i.south) -- (g\nm\i.south |- 0,-10);
      \node[blkw, minimum width=14mm, minimum height=6mm, inner sep=1pt, font=\figfont\footnotesize] (cpu\nm) at (-12,-17) {CPU};
      \node[blkw, minimum width=14mm, minimum height=6mm, inner sep=1pt, font=\figfont\footnotesize] (nic\nm) at (12,-17) {网卡};
      \draw[figline] (cpu\nm.north) -- (cpu\nm.north |- 0,-10);
      \draw[figline] (nic\nm.north) -- (nic\nm.north |- 0,-10);
    \end{scope}
  }
  % IB 交换机
  \node[blkw, minimum width=26mm, minimum height=6mm, font=\figfont\footnotesize] (ib) at (0,-32) {InfiniBand 交换机};
  \draw[ink, line width=1.0pt, dashed] (nicA.south) |- (ib.west);
  \draw[ink, line width=1.0pt, dashed] (nicB.south) |- (ib.east);
  % 图例
  \begin{scope}[shift={(-66,-41)}]
    \draw[accent, line width=1.4pt] (0,0) -- (7,0); \node[note, anchor=west] at (8,0) {NVLink：机内 GPU 间，约 900 GB/s};
    \draw[figline, line width=0.6pt] (60,0) -- (67,0); \node[note, anchor=west] at (68,0) {PCIe：主机与设备间，约 64 GB/s};
    \draw[ink, line width=1.0pt, dashed] (0,-5) -- (7,-5); \node[note, anchor=west] at (8,-5) {InfiniBand + RDMA：跨节点，约 50 GB/s（400 Gb/s）每端口};
  \end{scope}
\end{tikzpicture}
\caption{单机内与跨节点的典型数据通路（示意）。机内 GPU 经 NVLink/NVSwitch 互联，GPU 与 CPU、网卡之间经 PCIe 连接，不同服务器之间经 InfiniBand 网络通信。带宽为 H100 一代硬件的数量级参考值，具体取决于配置}
\label{fig:interconnect}
\end{figure}


把这几个数字放在一起，就能看出一条规律：不同层级链路的带宽相差一个数量级以上，因此\emph{通信越频繁的工作，越应该放在越近的地方}。第\,\ref{subsec:parallel}\,小节会看到，张量并行每一层都要通信，应当留在 NVLink 相连的同一台机器内；数据并行、流水线并行通信较少，才适合跨越 InfiniBand 网络。也正因为二者位于不同的互联层级、解决不同的问题，NVLink 与 InfiniBand 不应混为一谈。

\begin{pitfall}
“RDMA 完全绕过操作系统”只是一种简化说法。实际使用 RDMA 仍需注册内存、加载驱动、管理权限与连接；更准确的表述是：RDMA \emph{在数据传输路径上}省去了常规的 CPU 数据拷贝和部分内核协议处理，从而降低延迟和 CPU 开销。
\end{pitfall}

\subsection{驱动与计算软件栈：CUDA、CANN 与 NeuWare}
\label{subsec:stack}

有了芯片、存储和互联，还差最后一环：硬件只能执行机器指令，而算法开发者写的是 \texttt{torch.matmul} 这样的高层算子，谁来完成两者之间的翻译？答案是一套软件，它负责把高层算子翻译成设备上可执行的\term{内核}（\en{kernel}），并管理设备内存、任务调度和多卡通信。

这套软件称为\term{计算软件栈}，通常包括驱动、运行时、编译器、算子库和集合通信库五部分。以一次矩阵乘法为例（图\,\ref{fig:stack}）：框架收到 \texttt{matmul} 调用后选择具体实现；算子库（或编译器生成的代码）提供针对该硬件优化的内核；运行时负责分配显存、启动内核和同步；驱动把命令提交给硬件执行。多卡训练时，集合通信库负责完成 All-Reduce 等跨设备数据交换。

\begin{figure}[htbp]
\centering
\begin{tikzpicture}[fig, x=1mm, y=1mm]
  \def\cw{27}  % 列宽
  \def\rh{9.5} % 行距
  % 列标题
  \foreach \c/\t in {0/NVIDIA GPU, 1/华为昇腾 NPU, 2/寒武纪 MLU}
    \node[font=\figfont\footnotesize\bfseries] at (\c*30,6) {\t};
  % 行
  \foreach \r/\lab/\a/\b/\cc in {
      0/深度学习框架/PyTorch 等/PyTorch 等/PyTorch 等,
      1/算子库与编译器/{cuBLAS、cuDNN}/{CANN 算子库}/CNNL,
      2/集合通信库/NCCL/HCCL/CNCL,
      3/运行时与驱动/{CUDA Runtime}/{AscendCL 运行时}/{CNRT、CNDrv},
      4/硬件/GPU/NPU/MLU} {
    \node[note, anchor=east] at (-17,-\r*\rh) {\lab};
    \foreach \c/\txt in {0/\a,1/\b,2/\cc} {
      \ifnum\r=0
        \node[blkw, minimum width=\cw mm, minimum height=7mm, inner sep=1pt, font=\figfont\footnotesize] at (\c*30,-\r*\rh) {\txt};
      \else\ifnum\r=4
        \node[blk, minimum width=\cw mm, minimum height=7mm, inner sep=1pt, font=\figfont\footnotesize] at (\c*30,-\r*\rh) {\txt};
      \else
        \node[blka, minimum width=\cw mm, minimum height=7mm, inner sep=1pt, font=\figfont\footnotesize, line width=0.5pt] at (\c*30,-\r*\rh) {\txt};
      \fi\fi
    }
  }
  % 软件栈括号
  \draw[decorate, decoration={brace, amplitude=3pt}, line width=0.6pt]
    (75,-0.5*\rh-0.3*\rh) -- (75,-3.5*\rh+0.1*\rh) node[midway, right=3pt, note, align=left] {计算\\软件栈};
  % 调用方向
  \draw[->, line width=0.7pt, ink] (-42,1) -- (-42,-38) node[midway, left, note, text width=4mm, align=center] {调\\用\\方\\向};
\end{tikzpicture}
\caption{计算软件栈的层次与三种平台的对应组件。一次矩阵乘法调用自上而下依次经过各层；不同平台的组件职责相似，但实现与行为并不完全等价}
\label{fig:stack}
\end{figure}


不同硬件平台各有自己的软件栈：
\begin{itemize}
  \item \term{CUDA} 是 NVIDIA 面向自家 GPU 的并行计算平台和编程模型，包括编程接口、编译器、运行时，以及 cuBLAS（线性代数）、cuDNN（深度学习算子）、NCCL（集合通信）等库\cite{nvidia2024cuda}。经过近二十年积累，CUDA 生态是当前大模型训练与推理最主要的底层基础。
  \item \term{CANN} 是华为昇腾 NPU 的异构计算架构，提供编译器、运行时、算子库和 HCCL 集合通信库，承担与 CUDA 类似的职责。
  \item \term{NeuWare} 是寒武纪 MLU 系列芯片的基础软件平台，同样包含编程接口、编译工具与基础算子库。
\end{itemize}

软件栈的成熟程度，决定了硬件峰值性能能够发挥出多少。同一个模型在不同平台上的性能、算子支持范围和调试方法之所以不同，根源往往就在这一层。可见，CANN 等软件栈并非“PyTorch 的一个插件”；即使接口名称相似，也不能想当然地认为所有算子的行为和精度完全一致，迁移模型时需要逐项验证。

\begin{keypoints}
\begin{itemize}[leftmargin=1.5em]
  \item 芯片提供算力，存储提供容量与带宽，互联技术连接多卡多机，软件栈把硬件能力开放给上层。
  \item 判断硬件是否合适，要同时看算力、显存容量、显存带宽、互联拓扑和软件生态，而非单一峰值指标。
  \item 并行策略应与互联拓扑相匹配：通信越频繁的切分，越应放在带宽越高的链路上。
\end{itemize}
\end{keypoints}
